An additional concern was an unwanted contribution from Drell-Yan processes.
As Drell-Yan can only produce opposite-sign muons, by taking the ratio of the yields of opposite-sign muons to same-sign, and comparing to Monte Carlo we can approximate the contribution.
As mentioned in the last section, we expect the $\Delta p/p$ and $\Delta p/p$ significance of our single signal muons centered around 0, we expect the $\Delta p/p$ pair significance, simply defined by
$\Delta p/p_{pair sig}= \sqrt{\Delta p/p_{\mu1_{sig}}^2+\Delta p/p_{\mu2_{sig}}^2}$, to consequently be close to 0.
Taking the distribution as shown in left-hand side of Figure~\ref{fig:Drellyan}, we examine three regions of $\Delta p/p$ pair significance:
$0.0<\Delta p/p_{pair sig}<2.5$, $2.5<\Delta p/p_{pair sig}<4.5$, $4.5<\Delta p/p_{pair sig}<7.0$, which we expect to correspond to nearly pure signal,
a mixture of signal and background, and predominantly background, respectively.
From each of these regions we fit their $\Delta\phi$ distributions to a Gaussian fit centered around $\pi$, allowing for a combinatorial pedestal, and the natural underlying modulation, described by equation:

\begin{eqnarray}
  f=A+Bcos(2\Delta\phi)+\frac{C}{2\sqrt{2\pi}\sigma}e^{\frac{(\phi-\pi)^2}{2\sigma^2}}
\end{eqnarray}

and ascertain the yields.
An example fit is shown in the right-hand side of Figure~\ref{fig:Drellyan}.

\begin{figure} \centering
\includegraphics[width=0.5\textwidth]%
{figures/DrellYan/dpoppairsigzones}%
\includegraphics[width=0.5\textwidth]%
{figures/DrellYan/examplefit_awayside}
\caption{
REMAKE
\label{fig:Drellyan}
}
\end{figure}

Shown in TABLE, we compare the ratio of the yields of opposite-sign muons to same-sign for the three regions as well as our Monte Carlo sample.
As we see from the Monte Carlo, we expect the ratio to be about 2-1.
In the signal region of the Pb+Pb data, we do see a modest increase in the ratio from what we expect likely due to a Drell-Yan contribution.
We don't expect it to be large enough to make a significant difference and rather include it as a systematic uncertainty?
Additionally, we can see that as move into the higher $\Delta p/p$ significance region where we expect more background to be present, we see the ratio lessen as anticipated with increased, random, background.